% ===== Джерела до розділу 7. Ключі з префіксом ch07: =====
% Позиції ядра (bib-core.tex) тут НЕ дублюються.

% --- Динамічний аналіз та пісочниці ---
\bibitem{ch07:egele} Egele M., Scholte T., Kirda E., Kruegel C. A Survey on Automated Dynamic Malware-Analysis Techniques and Tools. \emph{ACM Computing Surveys}. 2012. Vol.~44, №~2. Art.~6. P.~1--42. DOI: 10.1145/2089125.2089126.

\bibitem{ch07:willems} Willems C., Holz T., Freiling F. Toward Automated Dynamic Malware Analysis Using CWSandbox. \emph{IEEE Security \& Privacy}. 2007. Vol.~5, №~2. P.~32--39.

\bibitem{ch07:cuckoo} Cuckoo Foundation. Cuckoo Sandbox: Automated Malware Analysis --- Documentation. URL: \url{https://cuckoo.readthedocs.io/}.

\bibitem{ch07:bayer} Bayer U., Comparetti P. M., Hlauschek C., Kruegel C., Kirda E. Scalable, Behavior-Based Malware Clustering. \emph{Proceedings of the 16th Annual Network and Distributed System Security Symposium (NDSS)}. San Diego: ISOC, 2009.

% --- Аналіз пам'яті ---
\bibitem{ch07:ligh} Ligh M. H., Case A., Levy J., Walters A. The Art of Memory Forensics: Detecting Malware and Threats in Windows, Linux, and Mac Memory. Indianapolis: Wiley, 2014. 912~p.

\bibitem{ch07:volatility} The Volatility Foundation. The Volatility Framework: Volatile Memory Artifact Extraction Utility Framework. URL: \url{https://www.volatilityfoundation.org/}.

% --- Ухилення від аналізу ---
\bibitem{ch07:chen} Chen X., Andersen J., Mao Z. M., Bailey M., Nazario J. Towards an Understanding of Anti-Virtualization and Anti-Debugging Behavior in Modern Malware. \emph{IEEE International Conference on Dependable Systems and Networks (DSN)}. Anchorage: IEEE, 2008. P.~177--186.

\bibitem{ch07:lindorfer} Lindorfer M., Kolbitsch C., Comparetti P. M. Detecting Environment-Sensitive Malware. \emph{Recent Advances in Intrusion Detection (RAID 2011)}. LNCS Vol.~6961. Berlin: Springer, 2011. P.~338--357.

\bibitem{ch07:kirat} Kirat D., Vigna G., Kruegel C. BareCloud: Bare-metal Analysis-based Evasive Malware Detection. \emph{Proceedings of the 23rd USENIX Security Symposium}. USENIX Association, 2014. P.~287--301.

% --- Поведінкове виявлення ---
\bibitem{ch07:bailey} Bailey M., Oberheide J., Andersen J., Mao Z. M., Jahanian F., Nazario J. Automated Classification and Analysis of Internet Malware. \emph{Recent Advances in Intrusion Detection (RAID 2007)}. LNCS Vol.~4637. Berlin: Springer, 2007. P.~178--197.

% --- Українські джерела ---
\bibitem{ch07:dynamic-ua} Застосування динамічного аналізу та пісочниць для дослідження шкідливого програмного забезпечення. \emph{Кібербезпека: освіта, наука, техніка}. 2024. URL: \url{https://csecurity.kubg.edu.ua/index.php/journal/article/view/615}.

\bibitem{ch07:behavior-ua} Поведінкові методи виявлення шкідливого програмного забезпечення на основі аналізу системних викликів. \emph{Захист інформації}. 2023. Т.~25, №~2.
